\chapter{Referencial Teórico}

\section{Ensino de Computação, inteligência artificial e formação na Educação Básica}

O ensino de Computação na Educação Básica não se limita ao uso de ferramentas digitais. Ele envolve conhecimentos e práticas que permitem compreender o funcionamento de sistemas computacionais. Recursos cotidianos baseados em inteligência artificial aproximam essa discussão da experiência dos estudantes, mas essa proximidade não substitui o ensino dos princípios que orientam esses recursos.

A Resolução CNE/CEB nº 1/2022 estabelece normas para a Computação como complemento à BNCC e atribui aos sistemas de ensino a definição de parâmetros e abordagens para sua implementação \cite{brasil2022computacao}.

O ensino de Computação difere da informática na educação, em que o computador apoia a aprendizagem de conteúdos de diversas áreas e o foco recai sobre sua inserção pedagógica \cite{valente2002uso}. A plataforma deste projeto também será um recurso didático, mas o objeto central de aprendizagem não é operar a interface nem estudar um conteúdo externo com auxílio do computador.

A proposta focaliza conceitos e práticas da Computação relacionados ao funcionamento e aos limites de sistemas de inteligência artificial. Espera-se que os estudantes interpretem a participação de dados e regras na saída gerada, comparando a finalidade dos modelos ao antecipar e explicar cada modificação.

A Matemática apoia essa leitura com gráficos e funções, além de noções de distância, média, erro e probabilidade, sem substituir o foco computacional das atividades. A Computação na Educação Básica se estrutura pelos eixos Pensamento Computacional, Mundo Digital e Cultura Digital, conforme as Diretrizes da Sociedade Brasileira de Computação e a Resolução CNE/CEB nº 1/2022 \cite{sbc2018diretrizes,brasil2022computacao}. Integrar esses eixos evita reduzir a plataforma ao manuseio de controles técnicos. No Ensino Médio, a seleção de conteúdos segue as habilidades da norma que articulam análise de dados e modelos computacionais à inteligência artificial \cite{brasil2022computacao}.

O pensamento computacional participa da proposta na organização do problema a partir de seus dados e regras, longe de equivaler a qualquer uso do computador. \citeonline{wing2006computational} chama a atenção para práticas como abstração e decomposição na formulação de soluções computacionais. Neste estudo, essa aproximação aparece quando o estudante separa uma situação em dados disponíveis e regras do modelo. Essa mobilização apoia a leitura do conteúdo específico da sequência, afastando a pretensão de medir um desenvolvimento geral do pensamento computacional ou de atestar ganhos diretos em outras disciplinas.

No Mundo Digital, os estudantes analisam como modelos processam representações numéricas e executam regras para produzir respostas que podem orientar decisões automatizadas. A Cultura Digital aparece no exame do uso de sistemas de inteligência artificial presentes no cotidiano.

\citeonline{touretzky2019envisioning} defendem que a educação em inteligência artificial para a Educação Básica deve ajudar os estudantes a compreender ideias que orientam sistemas inteligentes e seus impactos. O documento da AI4K12 sobre \textit{machine learning} \cite{touretzky2023machine} situa a aprendizagem a partir de dados como uma das grandes ideias da inteligência artificial e inclui, para a faixa do Ensino Médio, a diferenciação entre aprendizagem supervisionada e não supervisionada. Essas referências não são um currículo pronto para a realidade brasileira, mas oferecem linguagem conceitual para selecionar conhecimentos que não dependam de programação avançada.

\citeonline{camada2020ensino} identificaram possibilidades de pesquisa e aplicação para o ensino de inteligência artificial na Educação Básica a partir de uma revisão sistemática. \citeonline{duraes2021ensino} relatam uma sequência didática sobre fundamentos de IA com estudantes do Ensino Médio integrado. Esses estudos mostram que o tema já vem sendo trabalhado na Educação Básica. Na elaboração desta proposta, a mediação docente precisa considerar tanto o repertório matemático da turma quanto o tempo e a infraestrutura disponíveis.

Neste projeto, a inteligência artificial permite discutir como sistemas computacionais dependem de decisões humanas de modelagem. A questão educacional não será se o estudante consegue usar uma ferramenta generativa para terminar uma tarefa escolar. O que se buscará verificar é se, em situações delimitadas, o estudante consegue explicar como alterações em parâmetros e regras afetam previsões e classificações, de modo que a proposta permaneça vinculada ao ensino de Computação.

\section{\textit{Machine learning} introdutório no Ensino Médio}

O \textit{machine learning} ocupa lugar específico entre os conteúdos de inteligência artificial porque permite examinar como computadores identificam padrões em dados e os utilizam para produzir respostas em tarefas definidas. Ao mapearem unidades instrucionais destinadas à escola, \citeonline{marques2020teaching} observaram que vários materiais apresentam apenas aspectos básicos de \textit{machine learning} ou deixam partes do funcionamento do modelo em nível abstrato.

A proposta dispensa o detalhamento técnico de cada algoritmo. O estudante precisa compreender o suficiente para reconhecer como dados e regras participam da produção dos resultados do modelo.

\citeonline{martins2023findings} reúnem evidências de que estudantes iniciantes do Ensino Médio conseguem compreender conceitos básicos de \textit{machine learning} e aplicá-los em tarefas com algoritmos. \citeonline{sanusi2023systematic}, em revisão sobre a Educação Básica, também identificam experiências em diferentes etapas, com presença expressiva do Ensino Médio. As pesquisas apontam lacunas que vão da integração curricular e formação docente até a avaliação das intervenções. Assim, a adequação à faixa etária depende menos da idade e mais do recorte escolhido, das representações adotadas e do suporte oferecido na aula.

Esses estudos afastam a ideia de que o Ensino Médio seja inadequado para discutir modelos simples baseados em dados, embora alertem que o simples uso de recursos digitais está longe de garantir compreensão. As intervenções analisadas variam bastante em duração e objetivos conforme os recursos adotados. Uma experiência conduzida em uma escola específica dificilmente vira receita automática para outra. A literatura sustenta uma proposta viável, cabendo à aplicação revelar seus resultados.

A sequência não deve começar pela expressão matemática de um algoritmo, pois isso pode afastar estudantes que ainda consolidam a leitura de gráficos e noções de proporcionalidade. A manipulação de controles e a repetição de nomes técnicos também podem ocultar o raciocínio que se pretende desenvolver. A introdução ao tema precisa partir de problemas reconhecíveis para que a turma examine os dados e suas representações antes de discutir a regra que gera a saída. Essa progressão cria uma porta de entrada compatível com o nível de escolaridade previsto, sem diminuir o conteúdo de Computação.

O vocabulário técnico não é a única dificuldade esperada. O estudante pode precisar localizar pontos nos dois eixos para comparar distâncias, além de interpretar como as cores e curvas do gráfico expressam a regra do modelo. A leitura da visualização precisa ser ensinada como parte da atividade, pois observar uma mudança não significa explicar a relação que a produziu.

Há também um obstáculo matemático. Noções como média e distância, ou mesmo função e probabilidade, nem sempre estão estabilizadas na turma. A sequência passa longe de resolver, em oito encontros, lacunas acumuladas na escolarização matemática. Ela retomará esses conceitos apenas para responder a perguntas concretas sobre o modelo, mantendo a Matemática presente em um nível de formalização acessível.

Essa retomada não transforma o percurso em revisão extensa de Matemática antes do conteúdo de Computação. A necessidade aparece dentro de cada problema, como na comparação entre previsão e valor observado na regressão linear ou na distância entre exemplos no \textit{k}-NN. A mediação deverá identificar quando uma retomada curta basta e quando a dificuldade impede que a turma prossiga com a interpretação planejada. A seção 3.3 detalha como a Matemática entra nas representações.

Uma interface visual pode reduzir algumas barreiras e, ao mesmo tempo, ocultar outras. Um controle que altera o número de vizinhos torna uma mudança visível, mas não explica, por si só, por que certos pontos passam a decidir uma categoria. O recurso visual precisa ser acompanhado por mediação orientada, como se discute na seção 3.3.

A plataforma trabalha regressão linear, regressão logística, \textit{k}-vizinhos mais próximos, árvore de decisão e \textit{k}-means. A seleção abrange desde a previsão de valores numéricos e a classificação até o agrupamento de dados não rotulados. Em cada módulo, o estudante parte do problema proposto para levantar uma hipótese antes do experimento, explicando depois como a mudança nos dados ou no parâmetro alterou o gráfico e o resultado.

Os modelos foram escolhidos porque mostram mecanismos distintos em representações examináveis por estudantes iniciantes. Na regressão linear, o foco recai sobre a linha de tendência e a relação entre previsão, resíduo e erro. Já a regressão logística separa probabilidade e classificação, pois variar o limiar muda a decisão sem alterar o cálculo de base. Enquanto o \textit{k}-NN classifica pelo confronto com vizinhos semelhantes e a árvore de decisão explicita condições sucessivas, o \textit{k}-means forma grupos mesmo na ausência de categorias prévias para cada ponto.

Esse recorte reúne duas formas distintas de aprendizagem a partir de dados. Nos quatro primeiros módulos, que cobrem das regressões ao \textit{k}-NN e à árvore de decisão, os exemplos já trazem respostas conhecidas para orientar a previsão ou a escolha de categoria. No \textit{k}-means, não há categoria fornecida previamente, e o objetivo é formar grupos a partir de uma medida de proximidade. \citeonline{touretzky2023machine} situam essa diferenciação entre aprendizagem supervisionada e não supervisionada como uma das ideias centrais para a educação em IA. Neste projeto, ela será abordada com exemplos sintéticos e visualizações.

Os cinco módulos não devem ser assimilados com a mesma velocidade ou pelo mesmo caminho. A regressão linear pode se aproximar de conteúdos matemáticos já conhecidos, enquanto o agrupamento sem rótulos pode exigir mais tempo para que os estudantes diferenciem descobrir grupos e separar categorias previamente definidas. A retomada no encontro seguinte e a comparação coletiva no sétimo encontro evitam tratar cada modelo como episódio isolado.

O nível introdutório não exige ocultar os limites dos modelos. Uma reta pode resumir a tendência dos dados observados e ser afetada por pontos atípicos. Previsões fora do intervalo observado exigem cautela, pois a tendência pode não se manter. Uma porcentagem representada na regressão logística não é certeza.

No \textit{k}-NN, por exemplo, a decisão varia conforme os exemplos guardados e a escala adotada, além do número de vizinhos escolhido. Em árvores, o excesso de regras pode falhar diante de casos novos, enquanto os grupos do \textit{k}-means mudam segundo a quantidade e a posição inicial dos centroides. Tratar esses limites na leitura dos gráficos evita tanto a promessa de precisão automática quanto um debate moralizante e abstrato. Ao modificar dados e parâmetros nos controles de cada modelo, o estudante observa a consequência visual e discute qual regra a produziu.

\section{Representações, visualização e experimentação guiada}

Ferramentas visuais permitem introduzir \textit{machine learning} a estudantes iniciantes. \citeonline{gressevonwangenheim2021visual} identificaram ferramentas que favorecem a construção interativa de modelos, embora várias se concentrem na classificação de imagens e em atividades curtas. O produto proposto articula dados sintéticos e tabelas da própria turma em tarefas de regressão e classificação, além de agrupamento. Essa diferença de escopo não resolve, por si só, as limitações apontadas pela literatura.

A interação ganha sentido pedagógico quando o estudante investiga uma relação. A visualização pode tornar visível algo que permaneceria abstrato em uma explicação verbal, mas também pode dar a impressão de que observar uma mudança basta para compreendê-la. Os elementos visuais precisam ter significado explícito, e a regra do modelo deve estar acessível.

Os dados sintéticos permitem controlar o que será comparado e produzir situações em que a modificação de um elemento gera uma consequência perceptível. Não servem como previsão confiável sobre pessoas reais. Os dados organizados pela própria turma no segundo encontro mostram de onde vêm os números que serão interpretados. O cenário didático leva a turma a identificar a mudança ocorrida e a explicar qual elemento do modelo produziu aquele efeito.

Articular tabelas e gráficos a fórmulas, textos curtos e controles interativos mobiliza múltiplas representações. \citeonline{ainsworth2006deft} argumenta que representações diferentes podem cumprir funções complementares durante a aprendizagem, mas exigem tarefas de interpretação e integração. A plataforma não deve apenas reunir recursos na mesma tela. Cada representação precisa responder à pergunta examinada pela turma.

Qualquer ajuste nos controles da tela, seja na inclinação da reta e no limiar de decisão, seja nos parâmetros de vizinhos, condições ou centroides, pede uma questão que leve o estudante a antecipar e ler a mudança. Na regressão linear, os pares da tabela ganham forma no plano cartesiano para que a reta indique tendência ou previsão. Na regressão logística, a curva de probabilidade é lida em conjunto com o limiar que define a classe. Já no \textit{k}-NN, a posição do ponto consultado permite medir a distância até os exemplos para acompanhar a votação dos vizinhos. Enquanto nas árvores o percurso visual reflete as condições testadas, no \textit{k}-means os centroides funcionam como referências móveis que reorganizam os grupos, longe de serem pontos decorativos. Em todos os módulos, a visualização vincula os dados e a regra do modelo à saída gerada.

A atividade adapta a estratégia de prever, observar e explicar proposta por \citeonline{white1992probing}. Antes de usar o controle, o estudante registra uma hipótese oral ou escrita e, na sequência, acompanha como o gráfico altera a medida calculada e a decisão para explicar se o resultado confirmou sua expectativa. O ciclo evita a tentativa e erro em que parâmetros são movimentados sem relação clara entre entrada e saída.

A previsão torna visível uma ideia que poderia permanecer implícita. Se um estudante afirma que aumentar o número de vizinhos não alterará uma classificação e encontra decisão diferente, o professor pode perguntar quais exemplos passaram a ser considerados e por que a votação se alterou. Quando o estudante prevê que mover um centroide não terá efeito e observa a reorganização dos pontos, a explicação pode retomar a noção de proximidade. A hipótese inicial não precisa estar correta para ter valor pedagógico.

A aplicação está prevista para o laboratório de informática da escola, com acesso à internet, sujeito à confirmação das condições do campo. Planeja-se que cada estudante use um computador do laboratório ou, conforme a disponibilidade confirmada, um Chromebook, e a realização das atividades nos encontros busca reduzir as dificuldades decorrentes da restrição de acesso fora da escola. Para estudantes com insegurança no manuseio do computador, os controles deverão ser legíveis e indicar o que pode ser modificado em cada momento. Para aqueles que operam a interface com facilidade, o desafio deixa de ser o manuseio rápido do controle e passa a ser a explicação. A mediação docente, com suporte individual quando necessário, ajuda a evitar que a técnica de uso se torne obstáculo ou substitua a questão conceitual.

O controle interativo não será tratado como decoração ou recompensa após uma explicação pronta. Ele permite comparar uma consequência ao manter alguns fatores estáveis e alterar um dado ou parâmetro. A aula guiada orienta, inicialmente, uma alteração por vez, pois mudanças simultâneas dificultam atribuir uma consequência a uma causa percebida.

Depois da aula guiada, a exploração livre poderá ampliar as comparações. O professor deverá retomar a finalidade do experimento e solicitar explicações que não se limitem a descrever elementos apresentados.

As representações matemáticas serão apresentadas em camadas. O percurso parte de uma situação concreta e do uso de tabelas ou gráficos na leitura dos dados, passando por cálculos simples até chegar à fórmula que resume a regra discutida. Assim, tanto a equação da reta e o limiar quanto o cálculo de distâncias, condições ou médias dos centroides entram para explicar mudanças visíveis na tela, evitando a cobrança de memorização desarticulada.

Apresentar a Matemática em camadas não retira precisão dos conceitos. A organização em camadas define o que o estudante precisa para uma compreensão inicial e o momento em que cada noção passa a fazer sentido. Antes de calcular a distância entre dois pontos, o estudante precisa reconhecer por que a proximidade importa no \textit{k}-NN, e antes de discutir o erro de uma reta, precisa comparar uma previsão com um valor observado.

Introduzida após essas experiências, a fórmula ganha sentido prático por nomear uma relação que o estudante já conseguiu examinar visualmente durante o experimento.

\section{Avaliação da compreensão durante a intervenção}

Avaliar o ensino de \textit{machine learning} na Educação Básica permanece uma questão em aberto. \citeonline{rauber2023assessing} levantaram avaliações que vão de questionários a tarefas baseadas em desempenho e destacaram a necessidade de abordagens mais abrangentes.

\citeonline{gressevonwangenheim2022proposal} apresentam uma rubrica baseada em desempenho para conceitos e práticas de \textit{machine learning}, mostrando a importância de converter expectativas de aprendizagem em critérios observáveis. Embora voltado à classificação de imagens, esse modelo pode ser adaptado ao recorte desta pesquisa. Reconhecer o nome de um algoritmo está longe de provar que o estudante compreende a tarefa resolvida ou a ligação entre dados e saída. Como um acerto numérico pode decorrer de simples chute na leitura do gráfico ou da repetição de falas docentes, os instrumentos pedirão justificativas breves que revelem o raciocínio adotado.

A compreensão vai além do acerto em questões de múltipla escolha. O diagnóstico inicial e o núcleo comparável do instrumento final solicitarão operações cognitivas equivalentes em extensão e dificuldade. Enquanto registros individuais no segundo, quarto e sexto encontros acompanharão as previsões e explicações da turma, as notas do professor-pesquisador situarão as dúvidas surgidas e as intervenções nas condições de uso.

Longe de selecionar estudantes ou medir déficits, o diagnóstico inicial levantará concepções prévias. Isso ocorre quando alguém atribui decisões à inteligência artificial ignorando dados e regras, ou quando confunde previsão com classificação e supõe categorias prontas em agrupamentos. No instrumento final, o primeiro bloco retomará essas dimensões em questões equivalentes. Já o segundo bloco trará uma situação inédita para examinar a transferência conceitual, evitando equiparar essa tarefa mais complexa ao levantamento feito no início.

\citeonline{black1998assessment} discutem a avaliação formativa como parte do ensino, pois as informações produzidas pelos estudantes podem orientar decisões docentes. As perguntas na plataforma servirão para dar devolutiva imediata e organizar a conversa em sala, sem coleta digital de respostas. A análise confrontará o diagnóstico e o instrumento final com registros individuais selecionados e anotações docentes.

Durante a exploração, o professor poderá usar as previsões para identificar dificuldades que peçam retomada, dispensando o registro de cada fala como dado de pesquisa. Caso vários estudantes tomem uma probabilidade de 70\% como garantia de ocorrência, a aula poderá pausar a manipulação para distinguir estimativa de certeza. Da mesma forma, eventuais confusões entre \textit{k}-NN e \textit{k}-means serão tratadas retomando a presença ou ausência de rótulos nos exemplos. Tais intervenções pertencem ao ensino, cabendo ao registro posterior contextualizar as produções.

A análise se apoia em quatro dimensões que vão do reconhecimento da origem dos dados e da tarefa até a antecipação de ajustes, passando pela leitura de representações e dos limites na decisão automatizada. Uma rubrica em três níveis separará as respostas sem relação justificável daquelas com relações parciais ou explicações consistentes com o caráter introdutório da sequência, organizando o corpus sem gerar um índice isolado de competência.

Uma situação sobre consumo fictício de energia, por exemplo, poderá exigir que o estudante diferencie temperatura e consumo, evitando tratar todos os valores apenas como dados genéricos. Outras questões pedirão que ele antecipe o que acontece ao mover pontos e centroides na tela, ou ao mudar parâmetros como o limiar, as condições e o número de vizinhos. Assim, a leitura de tabelas e gráficos nos modelos será avaliada pela justificativa dada ao acerto.

Há diferença entre escrever que ``mudou porque o programa mudou'' e explicar que a alteração de limiar passou a classificar como alerta valores antes abaixo da regra estabelecida. A rubrica registrará essa diferença por dimensão, acompanhada de justificativa curta, sem produzir rótulos fixos para caracterizar estudantes.

A apresentação dos resultados combinará sínteses por dimensão e tabelas de transição entre momentos comparáveis, apoiadas por frequências descritivas com seus denominadores e breves excertos anonimizados. A análise acompanhará o que muda ou permanece em cada participante, evitando ranquear os estudantes. Já o desafio de transferência e o questionário de percepção terão seções próprias, situando a discussão nas condições que delimitam a intervenção.

\section{Síntese do referencial e posicionamento da proposta}

A literatura e os documentos curriculares definem o posicionamento da proposta. Nela, o \textit{machine learning} entra como conteúdo introdutório de Computação, superando o simples treinamento no uso de aplicativos. A sequência recorta ideias fundamentais e limites observáveis nas tarefas, sem antecipar uma formação universitária em algoritmos. Tanto os dados da própria turma quanto os cenários sintéticos visam mostrar a origem dos números e o efeito de alterar parâmetros. Dessa forma, a visualização sustenta comparações nas quais o estudante parte de uma hipótese para observar o gráfico e explicar como a regra do modelo transforma os dados em resultado.

A análise buscará evidências em produções individuais contextualizadas, indo além da contagem de acertos sem prometer efeitos causais generalizáveis.

Embora a literatura já reúna ferramentas e experiências no ensino de IA e \textit{machine learning} na Educação Básica, os estudos cobram maior clareza sobre conteúdos e estratégias nos contextos de aplicação, ao lado de avaliações mais rigorosas (\MakeUppercase{MARTINS; GRESSE VON WANGENHEIM}, \citeyear{martins2023findings}; \MakeUppercase{RAUBER; GRESSE VON WANGENHEIM}, \citeyear{rauber2023assessing}; \MakeUppercase{SANUSI} et al., \citeyear{sanusi2023systematic}). Assim, este projeto contribui ao analisar, em uma turma de Ensino Médio do próprio pesquisador, uma sequência didática que articula a visualização de cinco modelos a instrumentos com critérios explícitos.

A qualidade da pesquisa dependerá do alinhamento metodológico entre o referencial e o produto desenvolvido. Para isso, os recursos do site refletirão objetivos de aprendizagem presentes nas atividades e nos instrumentos de coleta, garantindo que as conclusões respeitem estritamente o alcance das evidências produzidas.
